\writeSectionFrame{\presentationTitle}{Interpreter}
\begin{frame}{Steckbrief}
	\begin{itemize}
 		\item Deutscher Name: Interpreter
 		\item Englischer Name: Interpreter
 		\item Gruppe: Verhaltensmuster
 	\end{itemize}
\end{frame}

\begin{frame}{Allgemeines Klassendiagramm}
    \begin{tikzpicture}[scale=0.7, transform shape]
          % Interpreter Pattern Classes
          \umlclass[x=0, y=6]{Client}{}{}
          \umlclass[x=0, y=3]{AbstractExpression}{}{}
          \umlclass[x=6, y=0]{TerminalExpression}{\textit{interpret(context)}}{}
          \umlclass[x=6, y=3]{NonTerminalExpression}{\textit{interpret(context)}}{}
          \umlclass[x=6, y=6]{Context}{}{}
          
          % Relationships
          
          \umlinherit{TerminalExpression}{AbstractExpression}
          \umlinherit{NonTerminalExpression}{AbstractExpression}
          \umluniassoc{Client}{Context}
          \umluniassoc{Client}{AbstractExpression}
        \end{tikzpicture}
\end{frame}
\note{
	Das Interpreter-Muster beschreibt den Aufbau einer Grammatik, damit sie von einem Interpreter
	ausgewertet werden kann. Außerdem muss der Ausdruck geparst werden. Der Parser ist jedoch nicht
	Teil des Musters. Es geht bei dem Muster darum, einen abstrakten Syntax-Baum zu repräsentieren
	und diesen zu interpretieren.
}

\begin{frame}{Abstrakter Syntax-Baum (AST)}
	\begin{center}
 		\includegraphics[scale=0.2]{\BILDERPFAD/interpreter1.png}
 	\end{center}
\end{frame}

\begin{frame}{Beispiele}
	\begin{itemize}
 		\item Römische Zahlen umrechnen
 		\item Buchtexte durchsuchen
 		\item SQL-Simulation
 	\end{itemize}
\end{frame}

\begin{frame}{Römische Zahlen umrechnen - Sprache}
	\begin{itemize}
 		\item M => 1000
 		\item C => 100
 		\item CD => 400
 		\item D => 500
 		\item CM => 900
 	\end{itemize}
\end{frame}

\begin{frame}{Römische Zahlen umrechnen - Sprache}
	\begin{itemize}
 		\item X => 10
 		\item XL => 40
 		\item L => 50
 		\item XC = 90
 		\item I => 1
 		\item IV => 4
 		\item V = 5
 		\item IX = 9
 	\end{itemize}
\end{frame}

\begin{frame}{Klassendiagramm römische Zahlen}
    \begin{tikzpicture}[scale=0.7, transform shape]
          % Interpreter Pattern Classes
          \umlclass[x=0, y=6]{RomanNumberParser}{}{}
          \umlclass[x=0, y=3]{Expression}{}{}
          \umlclass[x=6, y=-2]{ThousandExpression}{
            }{}
          \umlclass[x=6, y=3]{HundredExpression}{
          }{}

          \umlclass[x=0, y=0]{TenExpression}{
          }{}

          \umlclass[x=6, y=1]{OneExpression}{
          }{}
          \umlclass[x=6, y=6]{RomanNumberContext}{}{}
          
          % Relationships
          
          \umlinherit{ThousandExpression}{Expression}
          \umlinherit{HundredExpression}{Expression}
          \umlinherit{TenExpression}{Expression}
          \umlinherit{OneExpression}{Expression}
          \umluniassoc{RomanNumberParser}{RomanNumberContext}
          \umluniassoc{RomanNumberParser}{Expression}
        \end{tikzpicture}
\end{frame}

\begin{frame}[fragile]{Expression}
	\begin{exampleblock}{Expression.java}
 		\begin{lstlisting}
public abstract class Expression {
	public abstract String one();
	public abstract String four();
	public abstract String five();
	public abstract String nine();
	public abstract int multiplicator();

	public abstract boolean isStarting(String value);
}

 		\end{lstlisting}
 	\end{exampleblock}
\end{frame}

\begin{frame}[fragile]{Expression}
	\begin{exampleblock}{Expression.java}
 		\begin{lstlisting}
public void interpret(RomanNumbersContext romanNumbersContext) {
    String romanNumber = romanNumbersContext.getRomanNumber();
    
    if (romanNumber.length() > 0) {
        romanNumber = romanNumbersContext.getRomanNumber();
        while (isStarting(romanNumber)) {
            if (romanNumber.startsWith(nine())) {
                romanNumbersContext.add(9 * multiplicator());
                romanNumbersContext.removeFirstTwoRomanDigits();
            } else if (romanNumber.startsWith(four())) {
                // ...
            } else if (romanNumber.startsWith(five())) {
               // ...
            } else if (romanNumber.startsWith(one())) {
                // ...
            }
            romanNumber = romanNumbersContext.getRomanNumber();
        }
    }
}
 		\end{lstlisting}
 	\end{exampleblock}
\end{frame}

\begin{frame}[fragile]{TerminalExpression}
	\begin{exampleblock}{HundredExpression.java}
 		\begin{lstlisting}
public class HundredExpression extends Expression {
	@Override
	public String one() {
		return "C";
	}

	@Override
	public int multiplicator() {
		return 100;
	}
 		\end{lstlisting}
 	\end{exampleblock}
\end{frame}

\begin{frame}[fragile]{TerminalExpression}
	\begin{exampleblock}{HundredExpression.java}
 		\begin{lstlisting}
@Override
public boolean isStarting(String value) {
    if (value.startsWith("C")) {
        return true;
    }
    
    if (value.startsWith("CD")) {
        return true;
    } 
    // ...
    
    return false;
}
 		\end{lstlisting}
 	\end{exampleblock}
\end{frame}

\begin{frame}[fragile]{Context}
	\begin{exampleblock}{RomanNumbersContext.java}
 		\begin{lstlisting}
public class RomanNumbersContext {
	private String romanNumber;
	private int decimal;
	
	public RomanNumbersContext(String romanNumber) {
		this.romanNumber = romanNumber;
	}
	
	public void removeFirstRomanDigit() {
		romanNumber = romanNumber.substring(1);
	}
	
	public void removeFirstTwoRomanDigits() {
		romanNumber = romanNumber.substring(2);
	}
 		\end{lstlisting}
 	\end{exampleblock}
\end{frame}

\begin{frame}[fragile]{Context}
	\begin{exampleblock}{RomanNumbersContext.java}
 		\begin{lstlisting}
public void add(int number) {
    decimal = decimal + number;
}

public int getDecimal() {
    return decimal;
}

public String getRomanNumber() {
    return romanNumber;
}
 		\end{lstlisting}
 	\end{exampleblock}
\end{frame}

\frameWithImageOnly{\BILDERPFAD/interpreter2.jpg}
\note{
    Als nächstes implementieren wir eine Suche für ein Bibliothekssystem.
}

\begin{frame}{Buchsuche - Sprache}
	\begin{itemize}
 		\item AND => Literal1 vorhanden und Literal2 vorhanden
 		\item OR => Literal1 vorhanden oder Literal2 vorhanden
 		\item Literal => String-Ausdruck
 	\end{itemize}
\end{frame}

\begin{frame}{Klassendiagramm Buchsuche}
    \begin{tikzpicture}[scale=0.7, transform shape]
          % Interpreter Pattern Classes
          \umlclass[x=0, y=6]{BookSearch}{}{}
          \umlclass[x=0, y=3]{BookExpression}{}{}
          \umlclass[x=6, y=0]{AndExpression}{
            }{}
          \umlclass[x=6, y=3]{OrExpression}{
          }{}

          \umlclass[x=0, y=0]{TerminalExpression}{
          }{}

          \umlclass[x=6, y=6]{BookContentSearchContext}{}{}
          
          % Relationships
          
          \umlinherit{AndExpression}{BookExpression}
          \umlinherit{OrExpression}{BookExpression}
          \umlinherit{TerminalExpression}{BookExpression}
          \umluniassoc{BookSearch}{BookContentSearchContext}
          \umluniassoc{RomanNumberParser}{Expression}
        \end{tikzpicture}
\end{frame}

\begin{frame}{AST Buchsuche}
	\begin{center}
 		\includegraphics[scale=0.35]{\BILDERPFAD/interpreter3.png}
 	\end{center}
\end{frame}

\begin{frame}[fragile]{Expression}
	\begin{exampleblock}{BookExpression.java}
 		\begin{lstlisting}
public interface BookExpression {
	List<Book> interpret();
}
 		\end{lstlisting}
 	\end{exampleblock}
\end{frame}

\begin{frame}[fragile]{NonTerminalExpression}
	\begin{exampleblock}{OrExpression.java}
 		\begin{lstlisting}
public class OrExpression implements Expression {
	private Expression expression1;
	private Expression expression2;

	@Override
	public List<Book> interpret(BookContentSearchContext 
            interpreterBookContext) {
		List<Book> books = expression1.interpret(
            interpreterBookContext);
		books.addAll(expression2.interpret(
            interpreterBookContext));
		
		return books;
	}
}
 		\end{lstlisting}
 	\end{exampleblock}
\end{frame}

\begin{frame}[fragile]{TerminalExpression}
	\begin{exampleblock}{TerminalExpression.java}
 		\begin{lstlisting}
public class TerminalExpression implements Expression {
	private String searchTerm;

	@Override
	public List<Book> interpret(BookContentSearchContext 
            interpreterBookContext) {
		List<Book> foundBooks = new ArrayList<>();
		List<Book> books = interpreterBookContext.getBooks();
		
		for (Book book: books)  {
			String content = book.getContent();
			if (content.contains(searchTerm)) {
				foundBooks.add(book);
			}
		}
		
		return foundBooks;
	}
}
 		\end{lstlisting}
 	\end{exampleblock}
\end{frame}

\frameWithImageOnly{\BILDERPFAD/interpreter4.jpg}
\note{
    Als letztes Beispiel entwickeln wir eine kleine Simulation der Abfragesprache SQL. 
}

\begin{frame}{Klassendiagramm Sql-Simulation}
    \begin{tikzpicture}[scale=0.7, transform shape]
          % Interpreter Pattern Classes
          \umlclass[x=0, y=6]{SqlMain}{}{}
          \umlclass[x=0, y=3]{SQLExpression}{}{}
          \umlclass[x=6, y=-2]{NumericValue}{
            }{}
          \umlclass[x=6, y=3]{AndExpression}{
          }{}

          \umlclass[x=10, y=5]{IsPositiveExpression}{
          }{}          

          \umlclass[x=0, y=0]{BetweenExpression}{
          }{}

          \umlclass[x=6, y=1]{MaxExpression}{
          }{}
          \umlclass[x=6, y=6]{SqlTableContext}{}{}
          
          % Relationships
          
          \umlinherit{NumericValue}{SQLExpression}
          \umlinherit{AndExpression}{SQLExpression}
          \umlinherit{BetweenExpression}{SQLExpression}
          \umlinherit{MaxExpression}{SQLExpression}
          \umlinherit{IsPositiveExpression}{SQLExpression}
          \umluniassoc{SqlMain}{SqlTableContext}
          \umluniassoc{SqlMain}{SQLExpression}
        \end{tikzpicture}
\end{frame}

\begin{frame}{AST Sql-Simulation}
	\begin{center}
 		\includegraphics[scale=0.25]{\BILDERPFAD/interpreter5.png}
 	\end{center}
\end{frame}

\begin{frame}[fragile]{Expression}
	\begin{exampleblock}{SqlExpression.java}
 		\begin{lstlisting}
public interface SqlExpression {
	List<NumericValue> interpret(SqlTableContext sqlTableContext);
}

 		\end{lstlisting}
 	\end{exampleblock}
\end{frame}

\begin{frame}[fragile]{NonTerminalExpression}
	\begin{exampleblock}{LessThanExpression.java}
 		\begin{lstlisting}
public class LessThanExpression implements SqlExpression {
	private SqlExpression expression;
	
	@Override
	public List<NumericValue> interpret(SqlTableContext 
            sqlTableContext) {
		List<NumericValue> values = sqlTableContext.getValues();
		List<NumericValue> result = new ArrayList<>();
		NumericValue compareValue = expression.
            interpret(sqlTableContext).get(0);
		
		for (int i = 0; i < values.size(); i++) {
			// ...
		}
		
		return result;
	}
}
 		\end{lstlisting}
 	\end{exampleblock}
\end{frame}

\begin{frame}[fragile]{NonTerminalExpression}
	\begin{exampleblock}{BetweenExpression.java}
 		\begin{lstlisting}
public class BetweenExpression implements SqlExpression {
	private SqlExpression expression1;
	private SqlExpression expression2;

	@Override
	public List<NumericValue> interpret(SqlTableContext 
            sqlTableContext) {
		List<NumericValue> values = sqlTableContext.getValues();
		List<NumericValue> result = new ArrayList<>();
		
		for (int i = 1; i < values.size(); i++) {
			// ...
		}

		return result;
	}
}
 		\end{lstlisting}
 	\end{exampleblock}
\end{frame}

\begin{frame}[fragile]{TerminalExpression}
	\begin{exampleblock}{NumericValue.java}
 		\begin{lstlisting}
public class NumericValue implements SqlExpression {
	private int value;
	
	public int getValue() {
		return value;
	}

	@Override
	public List<NumericValue> interpret(SqlTableContext 
            sqlTableContext) {
		List<NumericValue> results = new ArrayList<>();
		results.add(this);
		return results;
	}
}
 		\end{lstlisting}
 	\end{exampleblock}
\end{frame}